% Рисунок с доски: характеристики уравнения переноса u_t + a u_x = 0 --
% параллельные прямые x - at = const с направляющим вектором (a, 1).
\begin{figure}[htbp]
\centering
\begin{tikzpicture}[>={Latex[length=4pt]}]
  \draw[->] (-0.3,0) -- (6.2,0) node[anchor=north east] {$x$};
  \draw[->] (0,-0.3) -- (0,3.3) node[anchor=south east] {$t$};

  % характеристики x = at + c, наклон dx/dt = a
  \foreach \c in {0.2, 1.2, 2.2, 3.2, 4.2} {
    \draw[thick, blue!60!black] (\c,0) -- ({\c+1.8},2.9);
  }

  % направляющий вектор
  \draw[->, red!65!black, very thick] (2.2,0) -- (2.82,1.0);
  \node[red!65!black, anchor=west, font=\footnotesize] at (2.88,0.92)
    {$\bar a = (a,\,1)$};

  \node[blue!60!black, anchor=north, font=\footnotesize] at (4.2,-0.06)
    {$x - at = \mathrm{const}$};
\end{tikzpicture}
\caption{Характеристики уравнения переноса $u_t + a u_x = 0$ в плоскости
$(x,t)$. Вдоль каждой прямой $x - at = \mathrm{const}$ концентрация не меняется:
профиль просто сносится средой со скоростью $a$.}
\label{fig:transport-lines}
\end{figure}
